% META: {"id":"1NSI-TC-EX-018","chapitre":"1NSI-TYPES-CONSTRUITS","type_objet":"exercice","capacites":["1NSI-TYPES-CONSTRUITS-C2"],"parcours":2,"competences":["programmer"],"duree_min":10,"mode_creation":"ex_nihilo","sources_inspiration":[],"status":"approved","origin":"needs_review","status_history":["needs_review","audited","accepted","approved"],"release_acceptance":"RELEASE_OWNER_BATCH_ACCEPTANCE","acceptance_closure_digest":"sha256:9b3ccf9a81c5520fbb7e03b7b2d3e7bf2057a02834b6d908bf3be5f49f3b553f","acceptance_receipt_digest":"sha256:4fad86f0282e0fa4e0419cca1ad6379ae7af5a280c04a4a90880f90a1c044242"}
% BEGIN-VERIFY
% def indice_maximum(tab):
%     idx = 0
%     for i in range(1, len(tab)):
%         if tab[i] > tab[idx]:
%             idx = i
%     return idx
% assert indice_maximum([3, 7, 2, 9, 1]) == 3
% assert indice_maximum([10]) == 0
% assert indice_maximum([5, 5, 5]) == 0
% END-VERIFY
\begin{exercice}{1NSI-TC-EX-018}{2}{10}
Écrire une fonction \lstinline{indice_maximum(tab)} qui prend un tableau non vide
d'entiers et renvoie l'indice de son plus grand élément. En cas d'égalité, on
renvoie le plus petit indice.

\textbf{Exemples :}
\begin{console}
>>> indice_maximum([3, 7, 2, 9, 1])
3
>>> indice_maximum([5, 5, 5])
0
\end{console}
\end{exercice}
